\documentclass[11pt, a4paper]{article}

\usepackage[T1]{fontenc}
\usepackage[utf8]{inputenc}
\usepackage{tgpagella}
\usepackage[spanish, es-noshorthands]{babel}
\usepackage{amsmath, amssymb, bm}
\usepackage{tabularx}
\usepackage{booktabs}
\usepackage[most]{tcolorbox}
\usepackage[margin=2.5cm]{geometry}
\usepackage{ragged2e}
\usepackage{fancyhdr}
\usepackage{tikz}

\pagestyle{fancy}
\fancyhf{}
\setlength{\headheight}{16pt}
\lhead{\textbf{UCB Sede Tarija} | Ing. Mecatrónica}
\chead{}
\rhead{IMT-121 Circuitos Electrónicos I | Actividad Autónoma EC3}
\lfoot{Prof: M.Sc. Ing. Bernardo Quiroga Turdera}
\rfoot{Página \thepage}
\renewcommand{\headrulewidth}{0.4pt}

\definecolor{ucbblue}{RGB}{0, 51, 102}

\begin{document}

\begin{tcolorbox}[colback=ucbblue!5!white, colframe=ucbblue, title=\textbf{Actividad Autónoma Formativa: Puente a Transitorios (EC3)}]
    \centering \textbf{Preparación Conceptual para la Unidad 3 (Para alumnos regulares)}
\end{tcolorbox}

\noindent\textit{Instrucciones: Esta actividad no es calificada y está diseñada para los alumnos que no han sido convocados a la evaluación sustitutiva. Lea el material introductorio sobre la física del almacenamiento de energía en circuitos y responda las preguntas cualitativas a continuación.}

\section*{1. Lectura: Inercia Eléctrica y Almacenamiento de Energía}
A diferencia de un resistor (que disipa energía como calor inmediatamente), un \textbf{Capacitor} acumula carga en un campo eléctrico y un \textbf{Inductor} acumula energía en un campo magnético. 

La física de estos componentes obliga a que el circuito no pueda reaccionar en cero segundos. 
\begin{itemize}
    \item Un capacitor se opone a cambios bruscos de \textbf{tensión}. Si usted cierra un switch sobre un capacitor descargado, la tensión $v_c$ arrancará en 0V y subirá lentamente (Régimen Transitorio) hasta alcanzar un Estado Estacionario.
    \item Un inductor se opone a cambios bruscos de \textbf{corriente}. Si se corta abruptamente un cable con corriente circulante en un inductor, este reaccionará violentamente para mantenerla, creando arcos de plasma (chispas).
\end{itemize}

\section*{2. Predicción Cualitativa (Cuestionario)}
Responda lógicamente basándose en las premisas de lectura:

\begin{enumerate}
    \item ¿Puede el voltaje almacenado en un capacitor ideal saltar instantáneamente (ej. de 0V a 12V en $t=0$) bajo una fuente DC ordinaria si cerramos un interruptor? Explique.
    \vspace{2cm}
    \item ¿Puede la corriente en un circuito inductivo saltar de la nada instantáneamente?
    \vspace{2cm}
    \item En la Unidad EC2 usábamos la Ley de Ohm ($V = IR$, ecuaciones algebraicas simples). En EC3, la inercia implica que la corriente de un capacitor es proporcional a la "velocidad de cambio" (derivada temporal) de su voltaje: $i_C(t) = C \frac{dv_C}{dt}$. \\
    Si aplicamos la Ley de Kirchhoff de Corrientes (LKC) en un nodo con capacitores y resistores, ¿Qué tipo de ecuaciones matemáticas inevitablemente aparecerán y tendremos que resolver?
    \vspace{2cm}
    \item En su cuaderno, intente bosquejar un gráfico $(x, y)$ donde el eje X sea el tiempo ($t$) y el eje Y sea el voltaje del capacitor $v_c(t)$ desde el momento $t=0$ en el que conectamos una batería hasta mucho tiempo después ($t \to \infty$).
\end{enumerate}

\end{document}
